% Generated draft: requires mathematical and source review.
\documentclass[11pt]{article}
\usepackage[margin=25mm]{geometry}
\usepackage{fontspec}
\setmainfont{Latin Modern Roman}
\usepackage{amsmath,amssymb,mathtools,unicode-math}
\setmathfont{Latin Modern Math}
\usepackage{enumitem,needspace,xurl,hyperref}
\setlength{\parindent}{0pt}
\setlength{\parskip}{5pt}
\setlength{\emergencystretch}{3em}
\newcommand{\N}{\mathbb{N}}
\newcommand{\Z}{\mathbb{Z}}
\newcommand{\Q}{\mathbb{Q}}
\newcommand{\R}{\mathbb{R}}
\newcommand{\C}{\mathbb{C}}
\newcommand{\F}{\mathbb{F}}
\newcommand{\A}{\mathbb{A}}
\newcommand{\PP}{\mathbb{P}}
\newcommand{\E}{\mathbb{E}}
\newcommand{\eps}{\varepsilon}
\newcommand{\dd}{\,\mathrm{d}}
\begin{document}
\section*{UnsolvedMath ID: 30002060}
\textbf{Source problem code:} catalogue entry
\subsection{Definitions and mathematical statement}
A finite geometric simplicial complex \(C\) in \(\R^m\) consists of straight simplices, contains every face of each simplex, and has the property that the intersection of two simplices is either empty or a common face. Its underlying space \(|C|\) is the union of its simplices. A linear triangulation of a convex \(d\)-ball means such a complex whose underlying space is convex and homeomorphic to the closed ball \(B^d=\{x\in\R^d:\|x\|\leq 1\}\). Here ``ball'' describes the topology of the underlying convex set; it does not require that set to be a round Euclidean ball.

An elementary simplicial collapse removes a pair of nonempty simplices \(\sigma,\tau\in C\), where \(\tau\) is a maximal simplex, \(\sigma\) is a codimension-one face of \(\tau\), and \(\tau\) is the only maximal simplex containing \(\sigma\). The resulting complex is \(C\setminus\{\sigma,\tau\}\). The complex \(C\) is collapsible if a finite sequence of elementary simplicial collapses reduces it to a single vertex; a complex consisting of one vertex is already collapsible.

Under the conventional finite-complex interpretation, the question is:
\[
\begin{gathered}
\text{For every pair of integers }0\leq d\leq m\text{ and every finite geometric simplicial complex }C\subset\R^m,\\
|C|\text{ convex and homeomorphic to }B^d
\quad\Longrightarrow\quad
C\text{ is collapsible}?
\end{gathered}
\]
The target is collapsibility of the given triangulation itself, without first subdividing it. No prescribed terminal vertex or requirement of collapse relative to the boundary is stated.

The supplied record does not explicitly specify finiteness or define ``convex ball.'' The formulation above makes the usual finite, closed-ball interpretation explicit; its agreement with the original report requires review. In particular, finite linear triangulations have polyhedral underlying spaces.

The supplied dated assessment attributes the question to Lickorish and reports it proved in dimensions below four and after at most one barycentric subdivision for linear subdivisions of polytopes, while reporting the assertion for every original linear triangulation in arbitrary dimension as open. These are supplied literature qualifications, not independently verified results or a solution asserted here.
\subsection{Short English statement}
Does every finite triangulation by straight simplices of a convex set homeomorphic to a closed ball collapse to one vertex by repeatedly deleting a simplex together with a free facet? The question concerns the original triangulation. Finiteness and the closed-ball interpretation are explicit conventions requiring source review. The supplied record labels the problem partially solved.
\subsection{Sources}
\begin{itemize}
\item Local UnsolvedMath registry, \texttt{data/unsolvedmath/problems.json}, permanent ID \(30002060\), source problem code \texttt{catalogue entry}, title: Collapsibility of Linear Triangulations of Convex Balls. This draft uses the supplied complete registry record.
\item Supplied original report citation: \emph{Triangulations (2013)}, Oberwolfach Reports, DOI \texttt{10.4171/owr/2012/24}, pp.\ 1405--1486. \url{https://doi.org/10.4171/owr/2012/24}. The record gives proposed year 2012; its citation title gives 2013. The report is described as documenting workshop discussions and open problems.
\item Supplied literature source labeled ``barycentric-subdivision theorem'': \url{https://doi.org/10.1007/s00454-019-00137-3}.
\item Supplied literature source labeled ``published theorem preprint'': \url{https://arxiv.org/abs/1709.07930}.
\item Literature assessment date: \texttt{2026-08-21}. This is supplied dated curation, not independently verified. The assessment attributes the question to Lickorish, reports results in dimensions below four and after at most one barycentric subdivision for linear subdivisions of polytopes, and labels the general original-triangulation question partially solved. The structured proposer field is empty.
\item The registry labels the statement \texttt{corrected\_verified} and reports a manual semantic, typographic, and LaTeX review of its cleanup. These are supplied editorial labels, not independent mathematical verification by this draft.
\end{itemize}
\end{document}
